\par\noindent\textbf{(1)} 群延迟每单位长度为 \(\beta_1=d\beta/d\omega=1/v_g\)。对波长求导 \(D=d\beta_1/d\lambda=(d\beta_1/d\omega)(d\omega/d\lambda)=-(2\pi c/\lambda^2)\beta_2\)。单位常写为 \(\mathrm{ps/(nm\cdot km)}\)，其符号与 \(\beta_2\) 相反。
\par\noindent\textbf{(2)} 若输入谱宽 \(\Delta\lambda\) 很小，群延迟差 \(\Delta\tau\approx L D(\lambda_0)\Delta\lambda\)。公式基于一阶波长展开而非把相速度差直接当群延迟；更宽谱应积分 \(L\int D(\lambda)d\lambda\)。
\par\noindent\textbf{(3)} 若 \(\beta_2(\omega_0)=0\)，二阶群速度色散消失，但 \(\beta_3=d^3\beta/d\omega^3\) 可能非零。对频率偏移 \(\Omega\)，群延迟 \(\tau_g(\Omega)=L[\beta_1+\beta_3\Omega^2/2+\cdots]\)，两侧对称谱的非线性时延仍引起波形畸变。
\par\noindent\textbf{(4)} 总色散为 \(D_{\rm total}=D_{\rm material}+D_{\rm waveguide}\)，其中波导项依模式横向场与芯包尺寸。可设计 \(D_{\rm total}=0\) 的波长，但这并不保证 \(\beta_3=0\) 或所有频带无色散，亦不自动消除自相位调制及偏振模色散。
